% =====================================================================
%  Branch and Price Solution Approach for 3D-SCPTOP
% =====================================================================
\section{Branch and Price Solution Approach}
\label{sec:math}

There are three structural difficulties with the 3D-SCPTOP model of
Section \ref{sec:model}. The first is that a priori we do not know how
many loading configurations are needed, so $|L|$ must be
overestimated; every configuration that turns out to be unnecessary
still carries its full block of placement variables. The second is
that the model replicates a complete two-dimensional packing problem
$|G|\cdot|L|$ times, and each replica contributes
$4\binom{|I|}{2}$ big-$M$ disjunctions
\eqref{sb:eq:nonov-x1}--\eqref{sb:eq:nonov-y2} whose linear
relaxation is weak. The third is that the model indexes its variables
by truck, and trucks are interchangeable: any two used trucks of the
same direction and period can swap their entire loads, so every
solution appears in the search a factorial number of times.

The third difficulty is the one that scales worst, because it does not
depend on the loading geometry at all. It is present in the volumetric
benchmark of Appendix \ref{appendix:A} exactly as in 3D-SCPTOP, and
Section \ref{sec:results} reports what it costs there: on the larger
instances the compact model's search tree grows until it reaches the
memory budget it is given, with its dual bound stalled. Operational
procurement planning lives at horizons of weeks over dozens of
trucks, so that is the regime the method has to serve, and it is the
regime in which a decomposition earns its place. A column generation
scheme based on branch and price is therefore developed.

The key components of the branch and price algorithm are the master
problem (MP), the restricted master problem (RMP), the pricing
problem (PP) and the branching policy. The MP is a reformulation of
the model that exploits its decomposable structure. Inspecting
Section \ref{sec:model}, every geometric constraint
\eqref{sb:eq:slot-active}--\eqref{sb:eq:util} carries a single
configuration index $l$, and every loading constraint
\eqref{sb:eq:slot-cap}--\eqref{sb:eq:stack-w-sup} carries a single
truck index $k$ and a single period $t$; the transport and inventory
block sees them \emph{only} through the aggregate flows $Q_{nt}$ and
$Q_{rt}$ of \eqref{sb:eq:flow-n}--\eqref{sb:eq:flow-r}. The model is
therefore block-angular, with one independent block per dispatched
truck. This suggests taking as the unit of decision neither the
individual RTI nor the loading configuration, but the
\emph{complete loaded truck}. A column is then a load, and the master
decides how many trucks carry it, so the trucks are never
distinguished in the first place and the symmetry disappears by
construction.

The MP is typically too large to be solved directly because it
involves a huge number of columns representing all possible truck
loads. To manage this complexity we consider only a smaller subset of
columns, which leads to the RMP. Solving the RMP provides an optimum
for the included subset, but not necessarily for the entire column
set. RMP integrality constraints are therefore relaxed to obtain dual
variables, and the PP uses those duals to identify columns of
negative reduced cost that should be added to the RMP. If no such
column exists, the integrality-relaxed MP has been solved to
optimality. If all the variables associated with the solution are
integers, an integer solution to the current branch has been found.
If not, the branch must be divided by adding further constraints.
Figure \ref{fig:B&P3D} exemplifies the algorithm.

\begin{figure}[h!]
    \centering
    \resizebox{0.8\columnwidth}{!}{%
\begin{tikzpicture}[node distance=2cm, auto, thick]
    \tikzstyle{block} = [rectangle, draw=white, fill=white, minimum height=2em,
                         text width=17em, text centered]
    \tikzstyle{line} = [draw, -Latex]

    \node [block] (MP) {\textbf{Master Problem}\\[2pt]
        \footnotesize all feasible truck loads $P^{p}\cup P^{s}$};
    \node[block, below=0.9cm of MP] (RMP) {\textbf{Restricted Master Problem}\\[2pt]
        \footnotesize loads generated so far, $\lambda\ge 0$};
    \node[block, below=0.9cm of RMP] (PP) {\textbf{Pricing Problem}\\[2pt]
        \footnotesize design one truck: 2D placement,\\
        \footnotesize layering and product assignment};

  \draw[bend right,left,->,align=left] (MP.west) to
        node {A reduced number \\ of columns} (RMP.west);
  \draw[bend right,left,->,align=left] (RMP.west) to
        node {Dual prices $\pi_{ct}$, $\nu_{gt}$} (PP.west);
  \draw[bend left,right,<-,align=left] (RMP.east) to
        node {If reduced cost is negative \\ add the truck load as a column; \\
              if not, branch} (PP.east);
  \draw[bend right,below,->,align=center] (PP.west) to
        node {Maximise $\sum_{c}\pi_{ct}a_{c}$ \\ subject to
              \eqref{eq:PPtype}--\eqref{eq:PPsymm}} (PP.east);
  \draw[bend left,below,<-,align=center] (RMP.west) to node {} (RMP.east);
\end{tikzpicture}
}
\caption{Branch and price scheme for 3D-SCPTOP. Every geometric,
height, crush-limit, payload and full-truckload constraint is enforced
inside the pricing problem, so each column is a physically loadable
truck.}
\label{fig:B&P3D}
\end{figure}


\subsection{Formulation of the Master Problem and the Restricted
Master Problem}
\label{sec:mp}

The key idea of the set-partitioning formulation for 3D-SCPTOP is to
enumerate all the feasible \emph{truck loads} with their associated
costs. A plant-bound truck load is a vector
$\mathbf{a}=(a_{n})_{n\in N}$ of full RTIs for which there exist a
slot count $m\le|I|$, stack types $h_{1},\dots,h_{m}\in H$,
orientations, bottom-left coordinates $(x_{i},y_{i})$, layer counts
$\lambda_{ir}$ and per-slot product loads $v_{in}$ such that
\begin{align}
&\text{the $m$ footprints, rotated as chosen, do not overlap and lie
   inside } pK\times qK,
   \label{eq:col-geom}\\
&\textstyle\sum_{r\in R(h_{i})} r_{pr}\lambda_{ir}\;\le\;rK
   &&\forall i\le m,
   \label{eq:col-height}\\
&\textstyle\sum_{n\in M(r)} v_{in}\;=\;n_{h_{i}r}\lambda_{ir}
   &&\forall i\le m,\;r\in R(h_{i}),
   \label{eq:col-assign}\\
&\textstyle\sum_{n\in N} w_{n}v_{in}\;\le\;W^{\max}
   &&\forall i\le m,
   \label{eq:col-crush}\\
&\textstyle\sum_{n\in N} w_{n}a_{n}\;\le\;pK^{w},
   \qquad
   \max\Bigl\{\tfrac{\sum_{n}\gamma_{p\rho(n)}a_{n}}{\Gamma_{K}},\;
              \tfrac{\sum_{n}w_{n}a_{n}}{pK^{w}}\Bigr\}\;\ge\;\alpha,
   \label{eq:col-truck}\\
&a_{n}=\textstyle\sum_{i\le m} v_{in}
   &&\forall n\in N.
   \label{eq:col-agg}
\end{align}
Let $P^{p}$ be the set of all such loads, and $P^{s}$ the analogous
set of supplier-bound loads $\mathbf{b}=(b_{r})_{r\in R}$, built with
folded heights $r_{sr}$, folded weights $w^{\mathrm{f}}_{r}$ and
folded volumes $\gamma_{sr}$; there no product assignment is needed,
because the weight of a folded RTI depends only on its type.

Two features of this definition matter. First, a column carries the
\emph{per-slot} product assignment $v_{in}$, so the crush limit
\eqref{eq:col-crush} is evaluated on the product mix actually riding
in that slot. The column therefore reproduces the period-exact cap
\eqref{sb:eq:stack-w-plant} of the compact model, and not the
conservative surrogate $\bar w_{r}=\max_{n\in M(r)}w_{n}$ that a pool
of capacity-only configurations is forced to use. Second, a column is
a feasibility certificate: the placement
$(h_{i},x_{i},y_{i},\text{rotation})$ witnessing
$\mathbf{a}\in P^{p}$ is retained, so any solution can be handed to a
loading crew as a drawing and re-verified independently of the solver
that produced it.

Constraint \eqref{eq:col-truck} is the same disjunction as
\eqref{sb:eq:util}--\eqref{sb:eq:utilw}, written on the load a truck
actually carries rather than on the capacity its configuration
provides. Both regimes occur in practice: coarse containers fill the
trailer by volume, dense ones by weight.

Let $\lambda^{g}_{pt}\in\mathbb{Z}_{\ge0}$ count the trucks
dispatched in direction $g$ and period $t$ carrying load $p$, and let
$a^{p}_{n}$ and $a^{p}_{r}$ be the RTIs of product $n$, respectively
of type $r$, that load $p$ contains. Every truck costs $cK$. The MP
of our problem is
\begin{align}
    \min\;\; cK \sum_{g\in G}\sum_{t\in T}\sum_{p\in P^{g}}\lambda^{g}_{pt}
    \;+\;\sum_{g\in G}\sum_{n\in N}\sum_{t\in T} ic_{gn} I_{gnt}
    \label{eq:MPObjective}
\end{align}
subject to
\begin{align}
    \sum_{p\in P^{p}} a^{p}_{n}\lambda^{p}_{pt} = Q_{nt}
      \quad & \forall n\in N,\;t\in T \label{eq:MPlinkp}\\
    \sum_{p\in P^{s}} a^{p}_{r}\lambda^{s}_{pt} = Q_{rt}
      \quad & \forall r\in R,\;t\in T \label{eq:MPlinks}\\
    \sum_{p\in P^{g}} \lambda^{g}_{pt} \le |K|
      \quad & \forall g\in G,\;t\in T \label{eq:MPfleet}\\
    \text{\eqref{sb:eq:inv-plant-main}--\eqref{sb:eq:safety}}
      \quad & \text{(inventory balances and coverage stock)}
      \nonumber\\
    \lambda^{g}_{pt}\in\mathbb{Z}_{\ge0}
      \quad & \forall g\in G,\;p\in P^{g},\;t\in T. \nonumber
\end{align}
Objective \eqref{eq:MPObjective} minimises transport plus inventory
holding cost. Constraints
\eqref{eq:MPlinkp}--\eqref{eq:MPlinks} state that the flow of full
and folded RTIs is exactly what the dispatched trucks carry, and
\eqref{eq:MPfleet} caps the fleet available per direction and period.
The coverage-stock rows \eqref{sb:eq:safety} are carried over for both
nodes and both kinds of stock, full and folded, exactly as the compact
model states them (Section \ref{sec:corrections}).

MP and 3D-SCPTOP have the same optimal value whenever $|L|$ is not
restricted a priori. Given a 3D-SCPTOP solution, each used truck
$(k,t)$ induces a load vector satisfying
\eqref{eq:col-geom}--\eqref{eq:col-agg}, obtained by reading the slot
data off its configuration and setting
$v_{in}=V^{\mathrm{slot}}_{kilnt}$; counting trucks by load vector
gives a MP solution of equal cost. Conversely, introducing one
configuration per distinct load used by a MP solution and assigning
trucks accordingly reconstructs a 3D-SCPTOP solution of the same
cost. This equivalence is what distinguishes the present approach
from a pool-based method such as PLC-SCPTOP: a predefined pool is a
\emph{restriction} and can only yield upper bounds of unknown
quality, whereas MP is the same problem written differently, so any
bound derived from it is a bound on the original.

The RMP replaces $P^{g}$ by a subset $\bar P^{g}$ and relaxes
integrality,
\begin{align}
    \min\;\; cK \sum_{g\in G}\sum_{t\in T}\sum_{p\in \bar P^{g}}\lambda^{g}_{pt}
    \;+\;\sum_{g\in G}\sum_{n\in N}\sum_{t\in T} ic_{gn} I_{gnt}
    \label{eq:RMPObjective}
\end{align}
subject to \eqref{eq:MPlinkp}--\eqref{eq:MPfleet} restricted to
$\bar P^{g}$, with $\lambda^{g}_{pt}\ge0$. Dual values are obtained
by solving the RMP. Let $\pi_{nt}$ and $\pi_{rt}$ be the duals of
\eqref{eq:MPlinkp}--\eqref{eq:MPlinks} and $\nu_{gt}\le0$ those of
\eqref{eq:MPfleet}. The reduced cost associated with a plant-bound
load $p\in P^{p}$ dispatched in period $t$ is
\begin{equation}
\bar c(p,t)\;=\;cK-\sum_{n\in N}\pi_{nt}a^{p}_{n}-\nu_{pt},
\label{eq:redcost}
\end{equation}
and analogously in the supplier direction. A load with negative
reduced cost can improve the current RMP solution. For each iteration
we solve the subproblem of finding the load of minimum reduced cost;
loads with negative reduced cost are added as new columns, the RMP is
re-solved to generate new duals, and the process repeats until no
further improving column is discovered.

Three implementation details of the RMP matter for what follows.
First, the RMP carries \emph{artificial} variables on the linking
rows \eqref{eq:MPlinkp}--\eqref{eq:MPlinks} and on the coverage-stock
rows, penalised at a large cost $M$ in the objective, so that it is
feasible for any pool and always yields duals; a solution with
positive artificials is not a plan. Second, it carries cumulative
cover cuts: by period $t$ the plant must have received cumulative
demand plus coverage stock, and dividing that requirement by the most
a single truck can carry bounds the number of trucks dispatched up to
$t-lt$ from below. They tighten the relaxation considerably, and their
duals, like those of every row the search later adds, are folded into
$\nu_{gt}$ so that \eqref{eq:redcost} remains the exact reduced cost.
Third, the RMP is \emph{period-scoped}: a column priced against the
duals of period $t$ receives a dispatch variable at $t$ alone, and only
the shared initial pool is offered to every period. This changes the
size of the restricted master, not the master --- the pricing
subproblem is still solved per $(g,t)$ over the whole of $P^{g}$, so a
sweep that finds nothing at any pair still proves that the restricted
LP equals the master LP.

\subsection{Lower bounds}
\label{sec:bounds}

Care is needed here, because the natural candidate is not a bound at
all. The RMP is a \emph{restriction} of the MP: it optimises over
$\bar P^{g}\subseteq P^{g}$, so its LP value satisfies
$z_{\mathrm{RMP}}\ge z_{\mathrm{MP}}$. Until column generation has
converged the RMP value is an \emph{upper} bound on the master LP,
and reporting it as a lower bound is simply wrong. Three rules decide
what may be reported; Figure~\ref{fig:bound} states them as a
decision.

\paragraph{A certificate needs a full sweep.}
When the PP has \emph{proven} that
$\zeta_{gt}:=\max_{p\in P^{g}}\sum_{c}\pi_{ct}a^{p}_{c}\le cK-\nu_{gt}$
for every $g$ and $t$, no column outside $\bar P^{g}$ has negative
reduced cost, so the current duals are feasible for the dual of the
full MP. Weak duality then gives
$z_{\mathrm{RMP}}=z_{\mathrm{MP}}\le z^{\star}_{\text{3D-SCPTOP}}$.
The proof is only as good as the exhaustiveness of the sweep behind
it. Two economies that are legitimate while columns are still coming
--- pricing only the pairs whose aggregates are fractional, and
skipping a pair whose duals round to something already priced ---
forfeit the right to certify, because a pair nobody looked at proves
nothing. Certification therefore runs one sweep over \emph{every}
$(g,t)$, with no memoisation and no restriction to the pairs a branch
has touched, and every call must come back \emph{proved}: an oracle
stopped on its time limit reports itself unproven, and a sweep with
one such pair proves nothing about the whole. When any of this fails
the node carries its parent's bound, which remains valid for the
child because the child's region is a subset of the parent's.

\paragraph{The bound is the raw LP value.}
The RMP objective includes the penalty $M$ on the artificial
variables. The cost of a plan is the value \emph{net} of that penalty,
and that is the right number to report as a cost; the bound is the
\emph{raw} value with the penalty included, because the penalty is
part of the LP whose dual the certificate rests on. Positive
artificials at convergence mean the region cannot be made feasible by
any column of $P^{g}$, so the region is infeasible in the true problem
and is pruned rather than bounded.

\paragraph{The restricted value is not a bound.}
At every node the search first solves the LP over the current pool.
That value sits \emph{above} the region's true LP value and is used
only as a filter: if even the over-estimate lies below the incumbent
the region cannot be pruned and the expensive sweep is not run. It is
never entered as a bound, and a search stopped by its deadline reports
the last \emph{certified} value it holds, not the last LP it solved.
Reduced costs oscillate here, and a sweep that finds nothing is
not convergence unless every pair was proved.

\paragraph{A bound by inclusion.}
The certificate above is expensive when the oracle is the constraint
programme of Section \ref{sec:pp}; the following observation makes a
cheaper one available for the geometric problem.

\begin{proposition}[bound by inclusion]
\label{prop:inclusion}
Let $P^{g}_{\mathrm{vol}}\supseteq P^{g}$ be the set of loads that fit
the trailer's volume and the payload and clear the full-truckload floor
on the binding resource, and let $z_{\mathrm{vol}}$ be the LP value of
the master \eqref{eq:MPObjective}--\eqref{eq:MPfleet} over
$P^{g}_{\mathrm{vol}}$, certified as above. Then
$z_{\mathrm{vol}}\le z^{\star}_{\text{3D-SCPTOP}}$, and the same holds
inside any region defined by bounds on the fleet vector.
\end{proposition}
\begin{proof}
A valid placement satisfying \eqref{eq:col-geom}--\eqref{eq:col-agg}
has total RTI volume at most $\Gamma_{K}$ and total weight at most
$pK^{w}$, and \eqref{eq:col-truck} is the same floor in both sets, so
$P^{g}\subseteq P^{g}_{\mathrm{vol}}$. The rows of the master --- the
linking, fleet, inventory and coverage rows, and any bound on
$n_{gt}$ or $X$ --- do not depend on the load model; only the column set
does. The 3D master is therefore the volumetric master with columns
removed, its LP value is at least $z_{\mathrm{vol}}$, and its optimum
is at least its LP value.
\end{proof}

The practical consequence is the part that matters. A geometric
pricing sweep costs $2|T|$ calls at the pricing limit and, on the
benchmark instances, does not certify inside a root window at all;
the volumetric relaxation --- a two-constraint knapsack oracle over the
same master --- certifies in a few seconds. The method therefore takes
its bound from the volumetric relaxation and asks geometric pricing
only for columns. It can run with \emph{no geometric pricing at the
root at all}: the plan comes from the seeded pool and the tree, the
bound comes from the relaxation, and the gap between them is a genuine
optimality gap for the 3D problem. This is a principled division of
labour rather than a shortcut, and Section \ref{sec:results} measures
what root pricing adds when it is given time.

\paragraph{Rounding the duals.}
The prices $\pi$ are irrational in general and must be rounded before
entering an integer-coefficient subproblem. We round them
\emph{upwards}, so the rounded objective dominates the true one and
the returned bound remains an over-estimate of $\zeta_{gt}$, as the
certificate requires. Rounding to nearest --- the obvious choice ---
makes the bound a slight under-estimate and silently invalidates the
certificate; with reduced costs at exactly zero it returns a value
just above the threshold and convergence is never declared.

\begin{figure}[h!]
\centering
\resizebox{\textwidth}{!}{%
\begin{tikzpicture}[>=Latex, font=\footnotesize, node distance=6mm,
  dec/.style={diamond, aspect=2.9, draw=black, fill=white, inner sep=1pt,
              align=center, font=\scriptsize},
  obox/.style={rectangle, draw=black, fill=white, align=center,
              inner sep=4pt, text width=42mm},
  why/.style={font=\scriptsize, align=left, text width=44mm, color=black!65},
  ln/.style={->, thick}]

  \node[dec] (d1) {was the sweep \emph{full}: every $(g,t)$,\\ no memoisation, no restriction\\ to the pairs the branch touched?};
  \node[dec, below=9mm of d1] (d2) {did \emph{every} pair come back\\ proved, none stopped on time?};
  \node[dec, below=9mm of d2] (d3) {are the artificial variables\\ zero at convergence?};
  \node[obox, below=9mm of d3] (rep) {\textbf{report the raw LP value}\\ penalty included; capped at the incumbent};
  \node[obox, right=22mm of d3] (inf) {\textbf{region infeasible}\\ prune; report nothing};
  \node[obox, left=18mm of d2] (parent) {\textbf{the node carries its}\\ \textbf{parent's bound}\\ valid, since the child's region\\ is a subset of the parent's};

  \draw[ln] (d1) -- node[right, font=\scriptsize] {yes} (d2);
  \draw[ln] (d2) -- node[right, font=\scriptsize] {yes} (d3);
  \draw[ln] (d3) -- node[right, font=\scriptsize] {yes} (rep);
  \draw[ln] (d3) -- node[above, font=\scriptsize] {no} (inf);
  \draw[ln] (d1.west) -- ++(-8mm,0) |- node[above, font=\scriptsize, pos=0.75] {no} (parent.north);
  \draw[ln] (d2.west) -- node[above, font=\scriptsize] {no} (parent.east);

  \node[why, right=6mm of d1.east, anchor=west]
      {a pair nobody looked at proves nothing; the economies that are
       sound while columns are still coming forfeit the certificate};
  \node[why, right=6mm of d2.east, anchor=west]
      {an oracle finding nothing is not a proof that nothing exists ---
       reduced costs oscillate, and a sweep stopped on time reports
       itself unproven};
  \node[why, left=6mm of d3.west, anchor=east]
      {the plan's cost is the value net of the penalty; the bound is the
       raw value, because the penalty is part of the LP the certificate
       rests on};
\end{tikzpicture}}
\caption{What may be reported as a lower bound. The restricted-master
LP value is an \emph{upper} bound on the master LP until pricing has
converged, so it is used only as a pruning filter; a node obtains a
bound of its own only from a full sweep in which every pair is proved,
and otherwise inherits its parent's. A quantity above a solution one
can exhibit is not a loose bound but a false statement, and it
inflates every optimality gap computed from it, so the reported value
is additionally capped at the incumbent.}
\label{fig:bound}
\end{figure}


\subsection{Pricing Subproblem}
\label{sec:pp}

We solve the pricing problem with a constraint programming approach.
The indices and parameters are the same as in the model proposed
before, but the decision variables do not depend on the configuration
index. Let $c$ range over products when $g=p$ and over RTI types when
$g=s$, let $\rho(c)$ be the RTI type carrying $c$, and let $S$ be the
slot bound \eqref{sb:eq:Ibound}. Table
\ref{tab:NomenclaturePricing3D} describes the new decision variables.

\begin{longtable}[hbt!]{|p{0.10\textwidth}|p{0.84\textwidth}|}
\caption{Nomenclature of the pricing subproblem}
\label{tab:NomenclaturePricing3D}\\
\hline
\multicolumn{2}{|l|}{\textbf{Decision variables}}\\ \hline
$u_{i}$        & Binary variable indicating whether slot $i$ is used in the new load\\
$\chi_{ih}$    & Binary variable indicating whether slot $i$ holds a stack of type $h$\\
$\varrho_{ih}$ & Binary variable indicating whether that stack is rotated by $90^{\circ}$\\
$\delta^{x}_{i},\delta^{y}_{i}$ & Footprint of slot $i$ along the $x$ and $y$ axes\\
$x_{i},y_{i}$  & Bottom-left coordinates of slot $i$ on the truck floor\\
$\ell_{ir}$    & Number of layers of RTI type $r$ in slot $i$\\
$v_{ic}$       & Number of RTIs of commodity $c$ loaded in slot $i$\\
$a_{c}$        & Number of RTIs of commodity $c$ in the new load\\
$\theta$       & Binary variable indicating whether the full-truckload requirement is met on volume rather than on weight\\
\hline
\end{longtable}

The objective is to find new truck loads with negative reduced cost.
Dropping the constant $cK-\nu_{gt}$ of \eqref{eq:redcost}, that is
\begin{align}
    \max z = \sum_{c}\pi_{ct}\,a_{c}
    \label{eq:PPObjective}
\end{align}
subject to
\begin{align}
 \sum_{h\in H}\chi_{ih}=u_{i},\qquad \varrho_{ih}\le\chi_{ih}
   \quad & \forall i\le S,\;h\in H \label{eq:PPtype}\\
 \delta^{x}_{i}=\sum_{h}\bigl(p_{h}\chi_{ih}+(q_{h}-p_{h})\varrho_{ih}\bigr)
   \quad & \forall i\le S \label{eq:PPdimx}\\
 \delta^{y}_{i}=\sum_{h}\bigl(q_{h}\chi_{ih}+(p_{h}-q_{h})\varrho_{ih}\bigr)
   \quad & \forall i\le S \label{eq:PPdimy}\\
 \textsc{NoOverlap2D}\Bigl(\bigl\{[x_{i},x_{i}+\delta^{x}_{i})\times
   [y_{i},y_{i}+\delta^{y}_{i})\bigr\}_{i\,:\,u_{i}=1}\Bigr)
   \quad & \label{eq:PPnooverlap}\\
 x_{i}+\delta^{x}_{i}\le pK,\qquad y_{i}+\delta^{y}_{i}\le qK
   \quad & \forall i\le S \label{eq:PPcontain}\\
 \sum_{i\le S}\sum_{h} p_{h}q_{h}\chi_{ih}\;\le\;pK\cdot qK
   \quad & \label{eq:PParea}\\
 \ell_{ir}\;\le\;\bigl\lfloor rK/r_{gr}\bigr\rfloor
   \sum_{h:r\in R(h)}\chi_{ih}
   \quad & \forall i\le S,\;r\in R \label{eq:PPlayercompat}\\
 \sum_{r\in R} r_{gr}\,\ell_{ir}\;\le\;rK
   \quad & \forall i\le S \label{eq:PPheight}\\
 \sum_{c\,:\,\rho(c)=r} v_{ic}=n_{hr}\ell_{ir}
   \quad & \forall i\le S,\;r\in R \label{eq:PPassign}\\
 \sum_{c} w_{c}\,v_{ic}\;\le\;W^{\max}
   \quad & \forall i\le S \label{eq:PPcrush}\\
 a_{c}=\sum_{i\le S} v_{ic}
   \quad & \forall c \label{eq:PPagg}\\
 \sum_{c} w_{c}\,a_{c}\;\le\;pK^{w}
   \quad & \label{eq:PPpayload}\\
 \sum_{c}\gamma_{g\rho(c)}a_{c}\;\ge\;\alpha\,\Gamma_{K}
   \quad & \text{if }\theta=1 \label{eq:PPftlvol}\\
 \sum_{c} w_{c}\,a_{c}\;\ge\;\alpha\,pK^{w}
   \quad & \text{if }\theta=0 \label{eq:PPftlwgt}\\
 \sum_{i\le S}\sum_{h} p_{h}q_{h}\chi_{ih}\;\ge\;\alpha\,pK\cdot qK
   \quad & \text{if }\theta=1 \label{eq:PPareacut}\\
 \sum_{h}h\,\chi_{ih}\;\ge\;\sum_{h}h\,\chi_{i+1,h}
   \quad & \forall i<S. \label{eq:PPsymm}
\end{align}
Constraints \eqref{eq:PPtype}--\eqref{eq:PPagg} restrict the loads
used in the pricing problem to feasible ones, reproducing
\eqref{eq:col-geom}--\eqref{eq:col-agg}.
Equation \eqref{eq:PParea} is a redundant area cut, and
\eqref{eq:PPareacut} its counterpart under the volume regime: a slot
contributes at most (footprint area)$\times rK$ of RTI volume, so the
utilisation floor implies a floor on the occupied floor area. The
latter is what allows the solver to find a \emph{first} feasible
layout quickly, which for tight instances is harder than optimising
once one is known. Constraint \eqref{eq:PPsymm} breaks the symmetry
among interchangeable slots by ordering them on the stack-type index.

The subproblem is stated as a constraint program because its hard
core is the disjunctive non-overlap relation
\eqref{eq:PPnooverlap}, which a CP solver propagates natively through
a global constraint, whereas a MILP must encode it with the four
big-$M$ disjunctions
\eqref{sb:eq:nonov-x1}--\eqref{sb:eq:nonov-y2} and accept the weak
relaxation they entail. Whichever technology solves it, the oracle
must solve \emph{this} problem: an earlier MILP oracle imposed the
full-truckload requirement as the sum of the two ratios in
\eqref{eq:col-truck} rather than as their maximum, solved a strictly
weaker problem, and returned a column that failed verification.

\paragraph{The volumetric oracle.}
The same master, with the same inventory balances, coverage stock and
fleet limit, is also solved over the loads of the volumetric benchmark
of Appendix \ref{appendix:A}: there a load is a column iff its volume
fits the trailer's, its weight the payload, and it clears the
full-truckload floor on the binding resource. Pricing is then a
two-constraint integer knapsack that closes in milliseconds. Running
both oracles through one master and one search is what makes the
comparison of Section \ref{sec:results} an experiment about the
loading representation rather than about two solution technologies,
and it is the setting in which the search of Section
\ref{sec:branching} can afford to certify at every node.

\paragraph{What a certificate costs.}
Figure~\ref{fig:pricing} shows the oracle and what it can and cannot
return. A certified node is a full sweep, $2|T|$ calls each closed to
proven optimality. On the volumetric oracle that is affordable at
every node. On the constraint programme it is not --- a single call on
the benchmark instances takes tens of seconds to minutes, and a sweep
is then the whole budget of a node --- so with the CP oracle the
certificate is taken from the volumetric relaxation
(Proposition~\ref{prop:inclusion}), at the root and in the one region
that can refute the incumbent outright (Section \ref{sec:branching}),
and geometric pricing is asked only for columns.

\begin{figure}[h!]
\centering
\begin{tikzpicture}[>=Latex, thick, node distance=7mm,
  box/.style={rectangle, rounded corners, draw=black!65, fill=black!4,
              text width=52mm, align=left, inner sep=5pt, font=\footnotesize},
  io/.style={rectangle, rounded corners, draw=black!45, fill=white,
             text width=28mm, align=center, inner sep=4pt, font=\scriptsize},
  side/.style={rectangle, draw=black!35, dashed, fill=white,
               text width=44mm, align=left, inner sep=4pt, font=\scriptsize}]

  \node[io] (in) {duals $\pi_{ct}$, $\nu_{gt}$\\[2pt] rounded \emph{upwards}};
  \node[box, right=9mm of in] (oracle)
    {\textbf{CP oracle: design one truck}\\[3pt]
     \textbullet\; \emph{stack types} --- one per distinct footprint;
        layers, heights, per-stack crush cap\\[2pt]
     \textbullet\; \emph{2D placement} --- slots on the floor, no
        overlap, orientation\\[2pt]
     \textbullet\; \emph{truck rules} --- payload, height, and the
        full-truckload floor};
  \node[io, right=9mm of oracle] (out)
    {a column: a fully specified, loadable truck};
  \draw[->] (in) -- (oracle);
  \draw[->] (oracle) -- (out);

  \node[side, below=8mm of oracle.south west, anchor=north west] (modes)
    {\textbf{what it returns}\quad $\max\sum_{c}\pi_{ct}a_{c}$\\
     a \emph{vertex} of $\operatorname{conv}(P^{g})$, for any duals\\[4pt]
     \textbf{what it cannot return}\\
     an interior load; those come from the construction
     (Section~\ref{sec:branching}) and join the same pool};
  \node[side, right=6mm of modes] (fast)
    {\textbf{what a certificate costs}\\
     one full sweep: every $(g,t)$, no memoisation, no restriction to
     the pairs a branch touched, every call closed to proven
     optimality; with this oracle that confines certification to the
     root and the region $X\le X^{\star}-1$};
  \draw[->, dashed, black!55] (modes.north) -- ++(0,6mm);
\end{tikzpicture}
\caption{The pricing subproblem. Geometry, height, crush limits,
payload and the full-truckload rule are enforced \emph{inside} the
oracle, so any column it returns is a physically loadable truck. It
maximises the dual-weighted load and therefore returns a vertex of
$\operatorname{conv}(P^{g})$ whatever the duals are; loads in the
interior of that hull (Figure~\ref{fig:hull}) are supplied by the
construction of Section~\ref{sec:branching}, not by pricing. A
certificate needs every call of a sweep closed to proven optimality,
which on this oracle is the whole budget of a node, and is therefore
spent only where a bound changes what the tree can conclude.}
\label{fig:pricing}
\end{figure}


\subsection{Initial Columns}
\label{sec:initcols}

The search starts from the plainest possible pool: one maximally full
truck per commodity and direction, obtained in the geometric case from
a constructive \emph{grid layout} --- tile the floor with the
commodity's stack type in its better orientation and fill each slot to
its height, weight and crush limit --- and in the volumetric case by
filling to capacity. No demand mixes, no diversification and no
reduced variants are seeded, so whatever the search then reaches, it
reached by pricing, branching and construction. The single-commodity
pool guarantees that no direction is ever left without columns, and
because these columns are shared they are offered to every period.

What such a pool, and any pool grown from it by pricing alone,
\emph{cannot} contain is shown in Figure~\ref{fig:hull} on a real
pool. For any dual vector, \eqref{eq:PPObjective} attains its
optimum at an extreme point of $\operatorname{conv}(P^{g})$, so
pricing returns extreme loads only. A balanced load lying inside the
hull is dominated componentwise by a fractional blend of
single-commodity trucks; the relaxation buys that blend for the price
of one truck, while an integer plan must pay for each truck it mixes.
No choice of duals returns the interior load, however the search moves
them. Interior loads are precisely what an integer plan needs when a
period's demand does not decompose into full trucks, and the method
therefore has a second generator for them, described next.

\begin{figure}[h!]
\centering
\begin{tikzpicture}[x=0.028cm, y=0.030cm, >=Latex,
  lab/.style={font=\scriptsize, align=left},
  lead/.style={black!55, thin, shorten >=2pt}]
  % axes
  \draw[->, thick] (-8,0) -- (262,0) node[below] {\footnotesize $a_{R_1}$};
  \draw[->, thick] (0,-8) -- (0,208) node[left] {\footnotesize $a_{R_2}$};
  \foreach \x in {50,100,150,200} \draw (\x,0) -- (\x,-5) node[below] {\tiny \x};
  \foreach \y in {50,100,150} \draw (0,\y) -- (-5,\y) node[left] {\tiny \y};

  % loads produced by the demand-derived price sweep (the seed pool)
  \foreach \p in {(0,127),(0,145),(0,163),(0,168),(0,182),(7,168),(14,168),
                  (16,100),(19,117),(21,168),(22,134),(25,151),(28,126),
                  (28,168),(142,0),(166,0),(190,0),(214,0),(238,0)}
      \fill[black!55] \p circle[radius=2.6];

  % upper-right frontier of the convex hull
  \draw[black!75, dashed, thick] (0,182) -- (28,168) -- (238,0);
  \node[lab, black!75, rotate=-31] at (168,64) {frontier of $\operatorname{conv}(P^{s})$};

  % dominating blend, and the load an integer plan needs
  \fill[blue!70!black] (74,131.2) circle[radius=3];
  \node[star, star points=5, star point ratio=2.2, fill=red!75!black,
        inner sep=1.7pt] at (74,126) {};
  \draw[blue!70!black, <->, thin] (74,129) -- (74,128.2);

  \node[lab, blue!70!black, anchor=west, text width=44mm] (bl) at (98,180)
      {$0.78\,(28,168)+0.22\,(238,0)=(74,131.2)$:\\ the blend the relaxation buys};
  \draw[lead, blue!70!black] (bl.south west) -- (78,133);

  \node[lab, red!75!black, anchor=west, text width=38mm] (tg) at (100,108)
      {$(74,126)$: the load a twelve-truck plan needs};
  \draw[lead, red!75!black] (tg.west) -- (79,125);

  % a dual vector, drawn inside the hull where there is room
  \draw[->, thick] (150,52) -- ++(26,20);
  \node[lab, anchor=north east] at (150,52) {$\pi$};
  \node[lab, anchor=west, text width=50mm] at (28,18)
      {for \emph{every} $\pi$, the maximum of $\pi^{\!\top}\!a$ is attained
       at a vertex of the frontier};
\end{tikzpicture}
\caption{Why the load a cheaper plan needs is never priced. Dots are
the supplier-bound loads obtained from a sweep of demand-derived price
vectors; the dashed line is the upper-right frontier of their convex
hull. The pricing objective is linear, so for \emph{any} dual vector
$\pi$ its optimum is attained at a vertex of that frontier. The load
$(74,126)$ lies strictly below it --- it is dominated componentwise by
the blend $0.78\,(28,168)+0.22\,(238,0)=(74,131.2)$ --- so no choice of
duals returns it. The relaxation buys that blend for one truck, while
an integer plan must pay for both of the trucks it mixes. The
construction of Section~\ref{sec:branching} reaches the interior
directly: it partitions a delivery vector into exact loads instead of
asking a linear objective for them.}
\label{fig:hull}
\end{figure}


\subsection{The fleet tree}
\label{sec:branching}

Branching on the master's own variables is the obvious rule and the
wrong one here. Bounding one $\lambda^{g}_{pt}$ removes one load
from the relaxation, but a Dantzig--Wolfe master of this kind holds
many columns that differ only in detail, so the relaxation answers by
shifting the same fractional weight onto a near-identical column. Run
to exhaustion on a benchmark instance with the pool frozen, that rule
branched thousands of times without ever producing an integral
relaxation. Branching on a per-period dispatch total has no symmetric
twin, but it has an escape of its own: the relaxation answers a bound
on one period by moving a truckload into a neighbouring period, at
the price of one period's holding cost, so the bound barely moves.

We therefore branch on the \emph{fleet}. Let
$n_{gt}:=\sum_{p\in P^{g}}\lambda^{g}_{pt}$ be the dispatch total of
direction $g$ in period $t$, $\mathbf{n}=(n_{gt})_{g,t}$ the fleet
vector, and $X:=\sum_{g,t}n_{gt}$ the total fleet. Both are integer in
every plan, and every bound placed on them is a row of the master that
already carries a dual, so the pricing subproblem is untouched. Given
an incumbent plan with fleet vector $\mathbf{n}^{\star}$ and total
$X^{\star}$, the tree has two levels that ask different questions;
Figure~\ref{fig:branching} draws them.

\paragraph{Level 1: the total fleet.}
Transport cost is exactly $cK\cdot X$, so a bound on the total has
nothing to compensate it. Three regions partition the space:
\begin{itemize}
\item $X\ge X^{\star}+1$ costs a whole truck with no offset. A plan in
  it improves on the incumbent only if the extra truck saves more than
  $cK$ minus the incumbent's gap in inventory, and the region's
  relaxation says at once whether that is possible; it is usually
  pruned by bound.
\item $X=X^{\star}$ fixes transport. Everything inside is inventory
  allocation, and this is where the second level goes.
\item $X\le X^{\star}-1$ is the only region that can beat the incumbent
  outright --- ``is there a plan with fewer trucks?'' --- and is
  therefore worth a certificate of its own.
\end{itemize}

\paragraph{Level 2: the fleet vector, inside $X=X^{\star}$.}
Fix an order on the pairs $(g,t)$ --- periods first, then direction.
The fleet vector of any plan in the region differs from
$\mathbf{n}^{\star}$ at a first index $j$, or nowhere. Branching at
the first pair where a solution may differ gives three children,
$n_{j}\le n^{\star}_{j}-1$, $n_{j}\ge n^{\star}_{j}+1$ and
$n_{j}=n^{\star}_{j}$, and the equality child is partitioned again
at index $j+1$ (Murty's scheme). The regions are disjoint and
exhaustive, and the equality child is expanded first so that the
neighbourhood of the incumbent is searched before its complement.
The partition is taken around the plan \emph{constructed at that
node}, not around the global incumbent, so a node that improves the
incumbent is immediately re-partitioned around the better plan and
nothing has to be restarted. The chain is cut at a fixed depth.

\paragraph{Processing a node.}
A node is a region imposed on the RMP as fleet rows. The LP over the
current pool is solved first; being an over-estimate of the region's
LP it is a filter and never a bound (Section \ref{sec:bounds}). If the
filter allows pruning and the certification policy admits the node,
column generation is run to a certificate; a certified value at or
above the incumbent prunes the node, an infeasible region is
discarded, and a certified value below the incumbent is recorded as
that region's bound. In every other case the node carries its
parent's bound. The construction below is then run \emph{inside} the
region --- it must stay inside, or the children it induces partition
around a point outside the node --- and if it returns a plan cheaper
than the incumbent the incumbent is replaced. Finally, if the region
fixes $X$, the Murty children are generated around the constructed
plan.

\paragraph{The construction.}
The primal heuristic is a flow-first construction in four steps. It
decides first how much moves and how many trucks carry it, and builds
the trucks afterwards.
\begin{enumerate}
\item \emph{Restricted master.} The LP over the pool, inside the
  region, gives the duals and the fractional point the node stands on.
\item \emph{Aggregate integer master.} On a fresh copy of the RMP the
  shipment quantities $Q_{nt}$, $Q_{rt}$ and the dispatch totals
  $n_{gt}$ are made integral, with the column variables left
  \emph{continuous}, and the region's bounds are imposed on the
  $n_{gt}$ as well as on the columns. This is where the objective is
  actually decided: it chooses the flows and the fleet, and it does so
  without a single oracle call. A solution carrying artificial slack
  beyond a small tolerance means the region is infeasible over this
  pool and the node is abandoned; a small slack means the integer
  solve ran out of time short of a clean solution, and the partitions
  built from it are still valid loads that step 4 may combine with the
  pool.
\item \emph{Partition.} For each direction and period, the delivery
  vector chosen in step 2 is split into \emph{exactly} $n_{gt}$ legal
  trucks by a small MIP: integer parts summing to the vector, each
  within the trailer's volume and payload and clearing the
  full-truckload floor on the binding resource, ordered by weight to
  break the symmetry among parts. If $n_{gt}$ parts cannot be formed,
  $n_{gt}+1$ is tried, provided the region allows it. A timeout is not
  a proof: only a model that comes back infeasible says the split is
  impossible.
\item \emph{Integer master over pool plus partitions.} The parts are
  added to the pool as columns and the integer master is solved over
  the enlarged pool, warm-started with the partition itself. Its
  solution is the node's plan, and its value the candidate incumbent.
\end{enumerate}

\paragraph{Two generators, one pool.}
Pricing returns extreme points of $\operatorname{conv}(P^{g})$; the
construction returns exact partitions of a delivery vector, which are
interior points that no dual vector reaches
(Figure~\ref{fig:hull}). A column is a feasible load, and a node's
constraints bound the multipliers $\lambda$, not the load, so a column
found anywhere is valid everywhere: every column from either generator
goes into the one shared pool, and a node that is pruned still leaves
its columns behind. That is the mechanism by which the tree improves
on the root: it is not that it reaches an integral relaxation --- it
essentially never does --- but that each region it visits asks the
construction for a plan under constraints the root never saw, and the
loads that plan needs are added for every later node to use.

\paragraph{What the tree reports.}
The incumbent is the cheapest verified plan found by any construction.
The root bound is the certificate obtained before any branching ---
under 3D loading, that of the volumetric relaxation
(Proposition~\ref{prop:inclusion}). The
global bound of a branch and bound is the minimum over its
\emph{frontier}, not its root: a region nobody expanded contributes
nothing, which is why a dive reports its root value however far it has
gone. At the end of the budget the regions still open are therefore
swept --- filtered by their pool LP and certified where the filter
allows --- and the global bound is the minimum over the certified
bounds of the open regions, floored at the root bound and capped at
the incumbent. Optimality is claimed \emph{only} when the root is
certified, no region remains open, no restart has broken the
partition, and the incumbent meets the global bound. None of the runs
reported in Section \ref{sec:results} meets these conditions, and the
paper claims no proof of optimality anywhere.

\begin{figure}[h!]
\centering
\resizebox{0.96\columnwidth}{!}{%
\begin{tikzpicture}[>=Latex, thick,
  lvl/.style={rectangle, rounded corners, draw=black!60, fill=black!3,
              text width=48mm, align=left, inner sep=4pt, font=\footnotesize},
  reg/.style={rectangle, draw=black, fill=white, inner sep=3pt,
              font=\scriptsize, align=center, minimum height=6mm},
  fix/.style={reg, fill=black!8},
  lab/.style={font=\tiny, align=center}]

  % --- the two levels, described ---
  \node[lvl] (l1) {\textbf{Level 1: the total fleet} $X=\sum_{g,t}n_{gt}$\\[2pt]
     transport cost is exactly $cK\cdot X$, so a bound on the total has
     nothing to compensate it --- unlike a per-period bound, which the
     relaxation answers by shipping in the neighbouring period};
  \node[lvl, below=6mm of l1] (l2) {\textbf{Level 2: the fleet vector, inside $X=X^{\star}$}\\[2pt]
     order the pairs $(g,t)$; at the first pair where a plan may differ
     from the incumbent branch $<$, $=$, $>$ and recurse on the equality
     (Murty). Every fleet vector differs from $\mathbf{n}^{\star}$ at a
     first index, so the regions are disjoint and exhaustive};
  \node[lvl, below=6mm of l2] (l3) {\textbf{At every node}\\[2pt]
     a plan is \emph{constructed} inside the region and its trucks join
     the shared pool; the children partition around that plan. The tree
     finds columns; the construction finds plans};

  % --- the tree itself ---
  \begin{scope}[xshift=70mm, yshift=4mm]
    \node[reg] (r) at (0,0) {root: incumbent $\mathbf{n}^{\star}$, $X^{\star}$};
    \node[reg] (a) at (-3.0,-1.5) {$X\ge X^{\star}+1$};
    \node[fix] (b) at (0,-1.5) {$X=X^{\star}$};
    \node[reg] (c) at (3.0,-1.5) {$X\le X^{\star}-1$};
    \draw[->] (r) -- (a); \draw[->] (r) -- (b); \draw[->] (r) -- (c);
    \node[lab, below=1pt of a, text width=26mm] {a whole truck with no offset:\\ usually pruned by bound};
    \node[lab, below=1pt of c, text width=26mm] {the only region that can beat\\ the incumbent outright: certified};

    \node[reg] (b1) at (-2.6,-3.4) {$n_{1}\le n^{\star}_{1}-1$};
    \node[fix] (b2) at (0,-3.4) {$n_{1}=n^{\star}_{1}$};
    \node[reg] (b3) at (2.6,-3.4) {$n_{1}\ge n^{\star}_{1}+1$};
    \draw[->] (b) -- (b1); \draw[->] (b) -- (b2); \draw[->] (b) -- (b3);

    \node[reg] (c1) at (-2.6,-5.0) {$n_{2}\le n^{\star}_{2}-1$};
    \node[fix] (c2) at (0,-5.0) {$n_{2}=n^{\star}_{2}$};
    \node[reg] (c3) at (2.6,-5.0) {$n_{2}\ge n^{\star}_{2}+1$};
    \draw[->] (b2) -- (c1); \draw[->] (b2) -- (c2); \draw[->] (b2) -- (c3);
    \node[lab, below=1pt of c2] {$\vdots$ \ to a fixed depth};
    \node[lab, anchor=west, text width=30mm] at (3.9,-4.2)
      {shaded: transport fixed;\\ everything inside is\\ inventory allocation};
  \end{scope}
\end{tikzpicture}}
\caption{The two levels of the fleet tree. Level 1 partitions on the
total fleet $X$, whose cost $cK\cdot X$ the relaxation cannot offset;
level 2 partitions the region $X=X^{\star}$ on the fleet vector itself,
one pair at a time around the plan constructed at each node. Bounds on
$X$ and on $n_{gt}$ are rows of the master that already carry duals,
so the pricing subproblem is untouched, and a column found in any
region is valid in every other.}
\label{fig:branching}
\end{figure}


\begin{figure}[h!]
\centering
\resizebox{\textwidth}{!}{%
\begin{tikzpicture}[>=Latex, font=\footnotesize, node distance=5mm,
  pbox/.style={rectangle, draw=black, fill=white, text width=62mm,
               align=center, inner sep=4pt, minimum height=8mm},
  dec/.style={diamond, aspect=2.8, draw=black, fill=white, inner sep=1pt,
              align=center, font=\scriptsize},
  note/.style={font=\scriptsize, align=left, text width=38mm, color=black!65},
  ln/.style={->, thick}]

  \node[pbox] (init) {one full single-commodity truck per commodity
        and direction};
  \node[pbox, below=of init] (cert) {\textbf{root certificate}: column
        generation to convergence --- full sweep, every $(g,t)$, every
        pair proved};
  \node[pbox, below=of cert] (cons) {\textbf{root construction}: aggregate
        integer master $\to$ partition $\to$ integer master over pool +
        parts};
  \node[pbox, below=of cons] (push) {push $X\ge X^{\star}+1$,\;
        $X\le X^{\star}-1$,\; $X=X^{\star}$};
  \node[pbox, below=of push] (pop) {pop a region; impose it as fleet rows;
        LP over the pool};
  \node[dec, below=6mm of pop] (filt) {pool LP $\ge$ incumbent\\ and policy certifies?};
  \node[pbox, below=7mm of filt] (cg) {certify the region; prune if
        $\mathrm{LB}\ge$ incumbent or infeasible};
  \node[pbox, below=of cg] (node) {construct \emph{inside} the region;
        its columns join the pool; update the incumbent};
  \node[dec, below=6mm of node] (fix) {region fixes $X$\\ and depth below cap?};
  \node[pbox, below=7mm of fix] (murty) {push the three Murty children
        around the plan just built};
  \node[pbox, below=of murty] (front) {\textbf{frontier}: filter and
        certify the regions left open; global bound $=$ min over them,
        floored at the root};
  \node[pbox, below=of front] (verify) {verify every dispatched truck
        against \eqref{eq:col-geom}--\eqref{eq:col-agg}};

  \draw[ln] (init) -- (cert);
  \draw[ln] (cert) -- (cons);
  \draw[ln] (cons) -- (push);
  \draw[ln] (push) -- (pop);
  \draw[ln] (pop) -- (filt);
  \draw[ln] (filt) -- node[right, font=\scriptsize, pos=0.4] {yes} (cg);
  \draw[ln] (filt.west) -- node[above, font=\scriptsize, pos=0.3] {no: parent's bound}
        ++(-14mm,0) |- (node.west);
  \draw[ln] (cg) -- (node);
  \draw[ln] (node) -- (fix);
  \draw[ln] (fix) -- node[right, font=\scriptsize, pos=0.4] {yes} (murty);
  \draw[ln] (fix.east) -- node[above, font=\scriptsize, pos=0.3] {no}
        ++(16mm,0) |- (pop.east);
  \draw[ln] (murty.east) -- ++(16mm,0) |- (pop.east);
  \draw[ln, dashed] (pop.west) -- ++(-14mm,0) |- node[left, font=\scriptsize, pos=0.25, align=right]
        {stack empty\\ or deadline} (front.west);
  \draw[ln] (front) -- (verify);

  \node[note, right=18mm of cert.east, anchor=west]
      {the RMP value is an \emph{upper} bound on the master LP until
       pricing converges; only a full, proved sweep turns it into a
       lower bound};
  \node[note, right=18mm of cons.east, anchor=west]
      {step 2 decides the objective with no oracle call; steps 3--4 build
       the trucks it implies};
  \node[note, right=18mm of pop.east, anchor=west]
      {a pool LP over-estimates the region's LP: it can say ``cannot
       prune'', never ``bound''};
  \node[note, right=18mm of node.east, anchor=west]
      {two generators, one pool: pricing gives vertices of
       $\operatorname{conv}(P^{g})$, the construction gives interior
       points; a pruned node still leaves its columns behind};
\end{tikzpicture}}
\caption{The fleet tree end to end. The search's value is not that it
reaches an integral relaxation --- it essentially never does --- but
that each region asks the construction for a plan under constraints the
root never saw, and the loads that plan needs are added to a pool every
later node can use. The bound is decided by certification alone: a
region that was not swept in full, with every pair proved, carries its
parent's bound, and the global bound is the minimum over the frontier
rather than the root value.}
\label{fig:procedure}
\end{figure}


\begin{algorithm}[!h]
\caption{The fleet tree for 3D-SCPTOP: branch and bound over the fleet
vector with the flow-first construction as primal}\label{alg:bp3d}
\begin{algorithmic}[1]
\State $\bar P \gets$ one full single-commodity truck per commodity and direction
       \Comment{Section \ref{sec:initcols}}
\State $\mathrm{LB}_{0} \gets \textsc{Certify}(\bar P,\ \text{root})$
       \Comment{full sweep, every $(g,t)$, every pair proved; else no bound.
                In 3D: on the volumetric relaxation, Prop.~\ref{prop:inclusion}}
\State $(\bar z,\mathbf{n}^{\star},\bar P) \gets \textsc{Construct}(\bar P,\ \text{root})$;
       \quad $X^{\star}\gets\sum_{g,t}n^{\star}_{gt}$
       \Comment{steps 1--4, Section \ref{sec:branching}}
\State $\mathcal{S} \gets \bigl[\,X\ge X^{\star}+1,\;\; X=X^{\star},\;\; X\le X^{\star}-1\,\bigr]$
       \Comment{a stack; the region that can beat the incumbent is expanded first}
\While{$\mathcal{S} \ne \emptyset$ and time remains}
   \State $\mathcal{R} \gets$ pop $\mathcal{S}$; impose $\mathcal{R}$ on the RMP as fleet rows
   \State $\tilde z \gets$ LP over $\bar P$ inside $\mathcal{R}$
          \Comment{over-estimates the region's LP: a filter, never a bound}
   \If{LP infeasible} \State \textbf{continue} \Comment{discarded} \EndIf
   \If{$\tilde z \ge \bar z$ and the policy certifies $\mathcal{R}$}
      \State $(\mathrm{LB}_{\mathcal{R}},\text{proved}) \gets \textsc{Certify}(\bar P,\mathcal{R})$
      \If{region infeasible or $\mathrm{LB}_{\mathcal{R}} \ge \bar z$}
         \State \textbf{continue} \Comment{pruned}
      \EndIf
      \State record $\mathrm{LB}_{\mathcal{R}}$ as an open bound
   \Else
      \State $\mathcal{R}$ carries its parent's bound
   \EndIf
   \State $(z,\mathbf{n},\bar P) \gets \textsc{Construct}(\bar P,\ \mathcal{R})$
          \Comment{inside the region; its columns join the shared pool}
   \If{no plan} \State price under the region's duals; \textsc{Construct} once more; \textbf{continue} if still none \EndIf
   \If{$z < \bar z$} \State $(\bar z,\mathbf{n}^{\star}) \gets (z,\mathbf{n})$
          \Comment{new incumbent} \EndIf
   \If{$\mathcal{R}$ does not fix $X$} \State push $\mathcal{R}\cap\{X\ge X_{2}+1\}$, $\mathcal{R}\cap\{X=X_{2}\}$, $\mathcal{R}\cap\{X\le X_{2}-1\}$ with $X_{2}=\sum_{g,t}n_{gt}$
          \Comment{level 1 recurses} \EndIf
   \If{$\mathcal{R}$ fixes $X$ and its depth $j$ is below the cap}
      \State $(g,t) \gets$ the $j$-th pair in the fixed order
      \State push $\mathcal{R}\cap\{n_{gt}\ge n_{gt}+1\}$,\;
             $\mathcal{R}\cap\{n_{gt}\le n_{gt}-1\}$,\;
             $\mathcal{R}\cap\{n_{gt}= n_{gt}\}$ onto $\mathcal{S}$
             \Comment{Murty, around the plan just constructed}
   \EndIf
   \State release $\mathcal{R}$
\EndWhile
\ForAll{$\mathcal{R}$ still on $\mathcal{S}$, while a reserve of time remains}
   \State filter by the pool LP; certify the survivors; record their bounds
   \Comment{the frontier}
\EndFor
\State $\mathrm{LB} \gets \max\bigl\{\mathrm{LB}_{0},\ \min\text{ over open and frontier bounds}\bigr\}$,
       capped at $\bar z$
\State verify every dispatched truck of the incumbent against
       \eqref{eq:col-geom}--\eqref{eq:col-agg}
\State \Return $(\bar z, \mathrm{LB})$;\ optimality is claimed only if $\mathrm{LB}_{0}$ was
       proved, $\mathcal{S}=\emptyset$, no restart occurred and $\bar z=\mathrm{LB}$
\end{algorithmic}
\end{algorithm}


\subsection{Two corrections to the compact model}
\label{sec:corrections}

Implementing 3D-SCPTOP surfaced two points worth recording, both
invisible at $lt=1$.

\paragraph{Inventory recursion under a lead time.}
Constraints \eqref{sb:eq:inv-plant-main} and
\eqref{sb:eq:inve-sup-main} carry inventory forward from period
$t-lt$, whereas inventory accumulates one period at a time and only
the \emph{arrival} is delayed. They must read
\begin{align}
I_{pnt}  &=I_{p,n,t-1}-d_{pnt}+Q_{n,t-lt}
    &&\forall n\in N,\; t\ge lt+1,
    \label{eq:fix-plant}\\
Ie_{str} &=Ie_{s,t-1,r}+Q_{r,t-lt}-\textstyle\sum_{n\in M(r)} d_{snt}
    &&\forall r\in R,\; t\ge lt+1.
    \label{eq:fix-sup}
\end{align}
With $lt=1$ the two versions coincide, which is why the defect
survived; with $lt\ge2$ the original silently discards $lt-1$ periods
of stock.

\paragraph{Coverage stock on folded RTIs.}
Constraint \eqref{sb:eq:safety} requires $Ie_{gtr}\ge ss_{gtr}$ at
\emph{both} nodes, and every model in this paper enforces it that way:
the master of Section \ref{sec:mp}, the volumetric benchmark of
Appendix \ref{appendix:A}, the company procedure PLC-SCPTOP and the
compact model itself. An earlier implementation imposed the folded
bound at the supplier only, on the argument that the plant has no
operational use for a stock of empties and that the extra row could
render an instance with a tight RTI fleet infeasible. The row is what
the model states, so it was reinstated, and the concern did not
materialise: on the largest instances of Section \ref{sec:results} the
root LP value is identical with and without it, no coverage row
carries artificial slack, and the row is tight in a minority of
periods without moving the optimum. Enforcing it everywhere also
removes an asymmetry: had the master enforced the plant side and the
comparison models not, the comparison would have been unfair in the
master's favour.

\subsection{The construction under 3D loading}
\label{sec:geo3d}

Step 3 of the construction builds trucks from a delivery vector with a
volumetric notion of a legal load. Over the volumetric oracle that
notion \emph{is} the definition of a column, and every part the
partition proposes is a column by construction. Under 3D loading a
load must also satisfy \eqref{eq:col-geom}--\eqref{eq:col-agg}. The
volumetric conditions are nevertheless \emph{necessary} for geometric
feasibility (Proposition~\ref{prop:inclusion}), so the partition MIP of
step 3 is kept as a \emph{relaxation} of the geometric partition and
each part it proposes is then tested for placeability. Step 2 makes no
oracle calls and is unchanged, so the objective is still decided before
any geometry is examined; what changes is only whether the trucks that
objective implies can be built.

\paragraph{Deciding placeability.}
A part is submitted to a placeability oracle that returns one of three
answers. A constructive \emph{shelf design} --- a small MIP that lays
the part's stacks in shelves along the trailer length, one integer
count per shelf, footprint and orientation --- supplies essentially
every positive answer, and its layout is passed through the
independent verifier of Section \ref{sec:results} before it counts; a
witness the verifier rejects is not a witness. The constraint programme
of Section \ref{sec:pp}, run in feasibility mode with the true slot
bound \eqref{sb:eq:Ibound}, supplies the \emph{infeasibility} proofs.
A part left undecided within the limit is neither accepted nor refuted;
retrying it longer does not help in practice, since such loads are hard
rather than slow. A placeable part becomes a column with its layout as
witness, exactly as a priced column does.

\paragraph{Two kinds of cut, and they are different claims.}
A part \emph{proved} unplaceable yields a \emph{dominance cut}
forbidding every load that dominates it componentwise. The cut is valid
because placeability is downward closed: removing an RTI from a
placeable layout leaves every position unchanged and every stack no
taller and no heavier, so a load dominating an unplaceable one is
itself unplaceable, and the cut removes the whole up-set of the refuted
load rather than a single point. The full-truckload floor is not part
of the placeability test, so it does not break the closure. A part left
\emph{undecided} is instead excluded as a single point, so that the
partition MIP proposes something else. That is a restriction of a
primal heuristic and never a validity claim: it cannot make the method
report an invalid plan, only fail to find one. The two constraints are
kept separate because they rest on different grounds, and only the
first may ever be read as knowledge about the instance.

\paragraph{Placeability caps.}
The maximum count of a single commodity in one truck is a scalar, and
by the same closure it is monotone, so it is bisected once per
direction and commodity at the start of a run --- a proved pricing
optimum tightens it, an undecided probe leaves the analytic cap in
force --- and handed to the partition MIP as an upper bound before it
solves. Without it a single-commodity dominance cut advances one unit
per round, from the volumetric fill down towards the placeable count,
at a placeability call and a partition re-solve each. The caps also
measure the volumetric relaxation's optimism directly: on one instance
it fills a truck with $272$ folded RTIs of a family of which $260$ can
be placed.

\paragraph{A stack relaxation inside the partition.}
The partition MIP additionally carries, per truck and footprint, an
integer number of stacks, with the crush limit and the height budget
imposed per stack and the stacks' floor areas summed against the
trailer floor. It is linear, valid --- every placeable load satisfies
it --- and it is exactly the constraint the volumetric relaxation
omits: a family whose RTIs are short enough for four layers by height
but heavy enough for only two by the crush limit is credited by volume
with twice its placeable count. The relaxation says nothing about the
two-dimensional arrangement of the stacks, which the placeability test
and the dominance cuts cover.

\paragraph{The pool and the bound.}
The starting pool is seeded without a solver: placeable pure trucks
and their reductions, and demand-shaped fills that the shelf design
places, every one of them verified. This is a requirement rather than
an optimisation, since over pure trucks alone step 2 is a hard
knapsack, and the volumetric fills of the single-commodity pool are
themselves not placeable in general. The root bound is the certified
value of the volumetric relaxation (Proposition~\ref{prop:inclusion});
geometric pricing, when given a share of the root budget, supplies
columns and nothing else. Section \ref{sec:results} reports what it
supplies.
